% !TeX root = Thesis.tex

\begin{lemma} \label{lemma::tetrahedron_subset_neighbourhood}
  Let $s_1,s_2,u,v$ form a $4$-clique in square $G$. If $N_G(s_1)\setminus\{s_1,s_2,u,v\}=N_G(s_2)\setminus\{s_1,s_2,u,v\}$, then
  \[
  N_G(u)\setminus\{s_1, s_2\}=N_G(v)\setminus\{s_1, s_2\}.
  \]
  That is, $N_G(u)\cup\{u\}$ and $N_G(v)\cup\{v\}$ induce the same clique structure in $G$.
\end{lemma}
\begin{proof}
  Let $x \in N_G(v)\setminus\{s_1,s_2\}$ be a vertex adjacent to $v$ in $G$ such that there exists a path of length at most two between $x$ and $v$ in $H$. If $x$ is adjacent to either $s_1$ or $s_2$, then by the assumption
  \[
  N_G(s_1)\setminus\{s_1,s_2,u,v\}
  =
  N_G(s_2)\setminus\{s_1,s_2,u,v\},
  \]
  where the vertices $s_1$ and $s_2$ contribute identically outside the clique. Hence any external connection from the clique through $v$ must also be realised through the same local structure containing $u$. Consequently $x$ is also at distance at most two from $u$ in $H$, and therefore
  \[
  x\in N_G(u)\setminus\{s_1,s_2\},
  \]
  which implies that $N_G(v)\setminus\{s_1,s_2\} \subseteq N_G(u)\setminus\{s_1,s_2\}$. By symmetry, interchanging the roles of $u$ and $v$ gives
  \[
  N_G(u)\setminus\{s_1,s_2\} \subseteq N_G(v)\setminus\{s_1,s_2\} 
  \]
  and hence  $N_G(u)\setminus\{s_1,s_2\} = N_G(v)\setminus\{s_1,s_2\}$. Therefore the vertices adjacent to $u$ and $v$ outside of $\{s_1,s_2\}$ have identical adjacency structure in $G$, and so $N_G(u)\cup\{u\}$ and $N_G(v)\cup\{v\}$ induce the same clique structure in $G$.
\end{proof}

\begin{proposition} \label{prop::textrahedron_decomposition_equality}
  Let $s_1,s_2,u,v$ form a $4$-clique in square $G$. If $\deg_G(u)=\deg_G(v)=d\geq3$, then one can choose to assign $x_{s_1u}=x_{x_2u}=x_{uv}=1$ in the tree root $H$ without loss of generality.
\end{proposition}
\begin{proof}
  
\end{proof}

\begin{proposition} \label{prop::tetrahedron_decomposition_feasibility}
  Let $s_1,s_2,u,v$ form a $4$-clique in square $G$. The linear program with Equation~\eqref{eq::tetrahedron_fix} never returns an infeasible solution for squares with tree roots. 
\end{proposition}
\begin{proof}
  
\end{proof}

\begin{proposition} \label{prop::tetrahedron_fixing}
  Let $H$ be a tree and let $u$ be an internal vertex with leaf children $s_1,\dots,s_k$ and parent $v$. Fixing any two leaves together with their internal vertex and parent uniquely determines the remaining leaf children of $u$ in the linear program.
\end{proposition}
\begin{proof}
  Since any tree may be rooted at an arbitrary vertex, assume without loss of generality that $H$ is rooted. Let $s_1,\dots,s_k$ be leaf children adjacent to an internal vertex $u$, and let $v$ be the parent of $u$. Consider any pair of leaves $s_i,s_j$, and suppose the internal branching vertex $u$ together with its parent $v$ have been fixed in the linear program. Now consider any remaining leaf $s_\ell$ adjacent to $u$. Since $s_\ell$ is a leaf in the root, its only neighbour in the tree is $u$. Furthermore, all leaf children of $u$ have identical neighbourhood structure in the square graph outside of $\{u,v,s_i,s_j\}$.

  Now consider any remaining leaf $s_\ell$ adjacent to $u$. Since every leaf adjacent to $u$ lies at distance two from $s_\ell$ through $u$, and $v$ lies at distance two through the path $s_\ell-u-v$, the neighbourhood of $s_\ell$ in the square graph is
  \[
  N_G(s_\ell) = \{u,v,s_1,\dots,s_{\ell-1},s_{\ell+1},\dots,s_k\}.
  \]

  Suppose for contradiction that $s_\ell$ is assigned to a vertex $w$ with $w\neq u$. Since $s_\ell$ is a leaf in the rooted tree, its neighbourhood in the square graph is determined by the vertices at distance at most two from $w$. Since $w\neq u$, the set of vertices at distance at most two from $s_\ell$ differs from that induced through $u$, giving
  \[
  N_G(s_\ell)
  \neq
  \{u,v,s_1,\dots,s_{\ell-1},s_{\ell+1},\dots,s_k\}.
  \]

  Hence either additional adjacencies are introduced into the square graph or existing adjacencies are removed and the neighbourhood of $s_\ell$ in the resulting square graph differs from that of the original graph, contradicting feasibility of the linear program. Thus the only feasible assignment is for $s_\ell$ to be adjacent to the already identified internal vertex $u$, implying that fixing any two leaves together with their internal vertex and parent uniquely determines all remaining leaf children adjacent to $u$.
\end{proof}